% !TeX root = ../presentacion.tex
% ===========================================================================
%  Datos de la portada, reparto de la exposición y metadatos del PDF.
%  Este es el único archivo que hay que editar con los datos del equipo.
% ===========================================================================

\newcommand{\tituloCharla}{Sistema de Reserva de Canchas Deportivas}
\newcommand{\subtituloCharla}{Microfrontends con Module Federation y microservicios Spring Boot}
\newcommand{\nombreCurso}{Desarrollo de Aplicaciones Empresariales}
\newcommand{\nombreCursoCorto}{Reserva de Canchas · Proyecto Integrador}
\newcommand{\nombrePrograma}{Maestría en Ingeniería de Software}
\newcommand{\nombreUniversidad}{\ph{Nombre de la Universidad}}
\newcommand{\tipoTrabajo}{Presentación final del proyecto integrador (\ent{6})}

% --- Integrantes del equipo ------------------------------------------------
% Mismos nombres que la portada del informe (informe/config/datos.tex).
\newcommand{\autorUno}{Pablo Jara Muñoz}
\newcommand{\autorDos}{Miguel Lara}
\newcommand{\autorTres}{Sebastián Uyaguari}
\newcommand{\autorCuatro}{Adrián Espinoza}

\newcommand{\nombreDocente}{\ph{Nombre del Docente}}
\newcommand{\ciudadFecha}{\ph{Ciudad}, \today}

% --- Reparto de la exposición ----------------------------------------------
% Quién expone cada bloque. Se imprime en la lámina separadora mientras la
% guía esté activa, de modo que al ensayar nadie dude de cuándo le toca.
% Si expone una sola persona, pon el mismo nombre en los cuatro.
\newcommand{\exponeArquitectura}{\ph{Integrante}}
\newcommand{\exponeFrontend}{\ph{Integrante}}
\newcommand{\exponeBackend}{\ph{Integrante}}
\newcommand{\exponeDemo}{\ph{Integrante}}

% --- Tiempo total ----------------------------------------------------------
% Ajusta este valor al que asigne el docente y revisa que los minutajes de
% los \bloque{...} sumen esta cifra menos el turno de preguntas.
\newcommand{\duracionTotal}{20 minutos + 5 de preguntas}

% --- Metadatos del PDF -----------------------------------------------------
% pdftitle y pdfauthor los fija beamer a partir de \title y \author de
% config/portada.tex; repetirlos aquí solo produce un aviso de hyperref.
\hypersetup{
    pdfsubject= {\nombreCurso\ --- \nombrePrograma},
    pdfkeywords = {microservicios, microfrontends, Module Federation,
                   Spring Boot, PostgreSQL, Docker}
}
